% Faithful conversion of ages-iii-sprint-3; source ODF IDs and text order retained.
\hypertarget{__RefHeading___Toc4821_3118176633}{}
\refstepcounter{subsection}\subsection*{4.3.4 Sprint 3}
\label{unit:ages-iii-sprint-3}
\addcontentsline{toc}{subsection}{4.3.4 Sprint 3}
% text-p-1237
\indent \originaltab{}Durante a Sprint 3, como eu e os demais AGES III já havíamos adiantado grande parte da infraestrutura, deploy em staging e esteiras de CI/CD nas \textit{sprints} anteriores, consegui focar mais diretamente no desenvolvimento de funcionalidades. Juntamente ao Gabriel Borges e à Isadora Spode (AGES I), desenvolvi as dependências da tela de pagamentos e, posteriormente, realizei a implementação da própria tela. Durante esse processo, revisei um dos PRs relacionados à funcionalidade e, após a aprovação, identifiquei um erro que exigiu um hotfix corretivo. Além disso, percebi a ausência de \textit{endpoints} responsáveis pela atualização do histórico de acessos após transações, entradas e saídas. Dessa forma, criei uma nova \textit{issue} para suprir essa necessidade futuramente.\par

% text-p-1238
\indent Também realizei, junto ao Erik Suris e ao Guilherme Khun (AGES III), a segregação dos ambientes de produção e \textit{staging} utilizando Terraform, criando configurações distintas para ambos os ambientes na AWS. Paralelamente, desenvolvi a tela de perfil e edição de perfil, mas durante a implementação encontrei inconsistências relacionadas aos campos de CPF e \textit{Cellphone}, que não estavam presentes nos \textit{models} da aplicação. Foi necessário interromper temporariamente a task para corrigir \textit{models}, DTOs, \textit{services} e criar uma nova \textit{migration} antes de retornar ao desenvolvimento da funcionalidade. Ainda nesse processo, o SonarCloud identificou uma vulnerabilidade potencial na lógica de validação de e-mails implementada no frontend, relacionada a problemas de \textit{backtracking} em regex. Após analisar o alerta, refatorei a implementação utilizando uma abordagem linear sem regex, reforçando ainda mais a importância das ferramentas automatizadas de análise estática dentro do desenvolvimento de software.\par

% text-p-1239
\indent Próximo à apresentação aos \textit{stakeholders}, o time novamente precisou realizar uma força-tarefa para revisar, corrigir e integrar funcionalidades pendentes na \textit{branch} principal. Eu e os demais AGES III nos dividimos nas revisões dos Pull Requests, e fiquei responsável pela análise de três PRs do \textit{frontend}. Durante esse processo, ficou evidente uma queda na qualidade geral das entregas da \textit{sprint}. Muitos PRs permaneciam abertos há dias com pipelines falhando e sem movimentação para correção, além da ausência de vinculação das issues aos Pull Requests, dificultando o processo de revisão e integração na main. Durante a retrospectiva da \textit{sprint}, também levantei junto ao time alguns pontos relacionados à falta de visão gerencial sobre o fluxo das \textit{tasks} no \textit{board}, principalmente pelo fato das \textit{dailies} terem funcionado muito mais como status \textit{reports} individuais do que como um acompanhamento efetivo do andamento do projeto.\par

\begin{reflection}
% text-p-1246
\indent \originaltab{}A Sprint 3 me fez perceber ainda mais o quanto uma pipeline de CI/CD bem estruturada é essencial para manter a qualidade de um projeto. Conforme a motivação do time foi diminuindo ao longo da \textit{sprint}, a qualidade das entregas e até das revisões também acabou sendo impactada. Nesse cenário, a pipeline teve um papel fundamental, identificando problemas, inconsistências e padrões que talvez passassem despercebidos pela correria da entrega. Algumas validações parecem pequenas ou até burocráticas no início, mas hoje vejo que elas ajudam a criar uma cultura de organização e responsabilidade no desenvolvimento. Foi algo que me marcou bastante e que quero levar para projetos futuros desde as primeiras \textit{sprints}.\par

% text-p-1247
\indent \originaltab{}Outro ponto que achei muito interessante foi perceber como ferramentas como o SonarCloud conseguem identificar vulnerabilidades que normalmente não pensamos durante o desenvolvimento. O alerta relacionado ao risco de \textit{backtracking} em regex me fez refletir sobre como aplicações reais precisam ser desenvolvidas pensando não apenas em funcionar, mas também em segurança, desempenho e confiabilidade. Mesmo sendo um ambiente acadêmico, a AGES proporciona experiências muito próximas do mercado, e aprender isso na prática acaba sendo extremamente valioso. Além disso, a retrospectiva da \textit{sprint} também foi importante para mostrar que a perda de ritmo impactou todo o time, não apenas no desenvolvimento das \textit{features}, mas também na gestão, revisões e organização do projeto. Ainda assim, foi positivo termos identificado esses pontos e alinhado uma Sprint 4 mais focada e consistente.\par

\end{reflection}
